\begin{figure*}[t]
\centering
\setcounter{subfigure}{0}
\begin{minipage}{145mm}
\centering
\ESPanel{.486\linewidth}{panels/fig11/a.pdf}{Composition summary and replacement design}%
\hfill%
\begin{minipage}[t]{.486\linewidth}
\vspace{0pt}\centering
\ESPanel{\linewidth}{panels/fig11/b.pdf}{Excess absolute change}
\par\vspace{3pt}
\ESPanel{\linewidth}{panels/fig11/c.pdf}{Signed difference}
\end{minipage}

\end{minipage}
\caption{Composition summaries and specified replacements reveal different aspects of member sensitivity. (a) Twelve-person teams on zsNo4HB9uLZ-16 combine six fixed members with six of twelve candidates. Horizontal labels give C1/C2/C3 member counts; points are mean Q, bars show fifth--95th percentiles, and $n$ gives the number of scorable teams. The member diagram below (a) defines the matched replacement used in (b--c). A separate 14-person A--B--C comparison shares seven retained members and exchanges seven-person blocks. (b) Excess magnitude $G_S$ compares cross-profile and within-profile Q changes; (c) signed $Q_C-Q_B$ shows whether C scores higher or lower than B. Each dot is an image mean; diamonds average the 69 image means from 146 scorable events (comparison definitions and supporting results in Supplements C17 and C19).}
\label{fig:composition-members}
\label{fig:matched-half}
\label{fig:team-comparison-design}
\end{figure*}
